
On considère la fonction à valeurs réelles $f(x)=4 e^{x / 4}-6$ dans l'intervalle $[0,4]$.
\begin{enumerate}
\item Montrer qu'il existe un zéro unique $\alpha$ pour la fonction $f$ dans l'intervalle $[0,4]$ et donner l'expression de $\alpha$.
\item Peut-on appliquer la méthode de la dichotomie pour calculer $\alpha$ ? Justifiez votre réponse.
\item Pour approcher le zéro $\alpha$ par la méthode du point fixe, on considère les suites récurrentes suivantes: $x_{(k+1)}=\phi_i\left(x_{(k)}\right), \mathrm{i}=1,2,3$, avec :

$$
\phi_1(x)=x+4 e^{x / 4}-6 \qquad \phi_2(x)=x-4+6 e^{-x / 4} \qquad \phi_3(x)=x-\frac{e^{x / 4}}{16}+\frac{3}{32}
$$


Etablir si chacune des suites converge vers $\alpha$ en utilisant la condition de convergence de la méthode du point fixe.
\end{enumerate}